\chapter*{Dear Programmer}
\markboth{Dear Programmer}{Dear Programmer}
\phantomsection
\addcontentsline{toc}{chapter}{Dear Programmer}

This book is for the kind of person who sees a blinking cursor and wonders
what it can do.

Maybe you have written programs before. Maybe this is the first time you have
typed a line of BASIC. Either way, you are in the right place. NovaBASIC is a
small language with direct access to a lively machine. You can print words,
draw pictures, make sounds, move sprites, save files, and build games one line
at a time.

The programs in this book are not magic tricks. They are meant to be taken
apart. Type them in, run them, change a number, break something, fix it, and
then make the game yours.

\begin{typeinbox}[title=The Most Important Rule]
Do not be afraid of mistakes. A typing error is not a disaster. It is the
computer telling you exactly where the next little mystery is.
\end{typeinbox}

Nova gives you more screen space, more memory, hardware sprites, colorful
graphics, useful sound, and fast helpers for moving data around. The listings
start simply because that is how good programs usually start. Once the simple
version works, you can make it faster, brighter, stranger, and more your own.

So mount the disk, turn to a listing, and type. The first run might not work.
The second or third one probably will. That feeling is worth chasing.

\vspace{1.5em}
\noindent\textbf{Welcome to Nova Fun \& Games.}
